% 使用 XeLaTeX 编译两次；resume.cls 保持原样。
\documentclass{resume}

\usepackage{tabularx}
\usepackage{needspace}
\usepackage{xurl}
\usepackage{lastpage}

\ResumeName{邹永康}

\hypersetup{
  pdftitle={邹永康 - Agent Harness 工程师},
  pdfauthor={邹永康}
}

\urlstyle{same}

% =========================================================
% 页面布局
% 字体沿用 resume.cls，不重新指定中英文字体
% =========================================================

\geometry{
  a4paper,
  left=15mm,
  right=15mm,
  top=13mm,
  bottom=14mm,
  footskip=7mm
}

\newcommand{\BodyFont}{%
  \zihao{5}
  \setlength{\baselineskip}{15pt}
}

\newcommand{\DetailFont}{%
  \zihao{-5}
}

\newlength{\EntryGap}
\setlength{\EntryGap}{9pt}

\newlength{\BulletGap}
\setlength{\BulletGap}{3pt}

\setlength{\parindent}{0pt}
\setlength{\parskip}{0pt}
\setlength{\tabcolsep}{0pt}

\raggedright
\raggedbottom

\emergencystretch=2em
\hyphenpenalty=10000
\exhyphenpenalty=10000
\widowpenalty=10000
\clubpenalty=10000

% 保留 cls 的章节字体和横线，仅调整间距
\ctexset{
  section={
    indent=-1em,
    beforeskip=10pt,
    afterskip=7pt
  }
}

% =========================================================
% 条目布局
% 保留 ResumeItem 接口与原有字号、斜体
% 第一行：公司 / 项目名称 + 右对齐日期
% 第二行：五号斜体职位 / 奖项
% =========================================================

\RenewDocumentCommand{\ResumeItem}{omoo}{%
  \par
  \Needspace{6\baselineskip}%

  \IfValueTF{#1}
    {\pdfbookmark[2]{#1}{subsec:\arabic{resumebookmark}}}
    {\pdfbookmark[2]{#2}{subsec:\arabic{resumebookmark}}}%

  \stepcounter{resumebookmark}%

  \begin{tabularx}{\linewidth}{
    @{}
    >{\raggedright\arraybackslash}X
    @{\hspace{12pt}}
    r
    @{}
  }
    {\zihao{-4}\heiti\bfseries #2}
      & {\DetailFont\IfValueT{#4}{#4}}\\[2pt]

    \multicolumn{2}{@{}p{\linewidth}@{}}{
      {\zihao{5}\textit{\IfValueT{#3}{#3}}}
    }\\
  \end{tabularx}\par\nobreak
}

\newenvironment{resumeBullets}{%
  \begin{itemize}[
    label=\textbullet,
    leftmargin=12pt,
    labelsep=5pt,
    itemsep=\BulletGap,
    parsep=0pt,
    topsep=3pt,
    partopsep=0pt
  ]
}{%
  \end{itemize}
}

\newcommand{\ProjectLink}[1]{%
  \par\nopagebreak
  {\DetailFont
    \textbf{链接:} \url{#1}\par
  }
}

% 教育经历：学校与日期一行，学位与成绩分别放在下一行
\newcommand{\Education}[4]{%
  \par
  \Needspace{4\baselineskip}%

  \begin{tabularx}{\linewidth}{
    @{}
    >{\raggedright\arraybackslash}X
    @{\hspace{10pt}}
    r
    @{}
  }
    {\zihao{5}\heiti\bfseries #1}
      & {\DetailFont #2}\\[2pt]

    \multicolumn{2}{@{}p{\linewidth}@{}}{
      {\zihao{5}\textit{#3}}
    }\\[1pt]

    \multicolumn{2}{@{}p{\linewidth}@{}}{
      {\DetailFont #4}
    }\\
  \end{tabularx}\par
}

% =========================================================
% 页脚
% =========================================================

\pagenumbering{arabic}

\pagestyle{fancy}
\fancyhf{}

\renewcommand{\headrulewidth}{0pt}
\renewcommand{\footrulewidth}{0pt}

\fancyfoot[L]{%
  \fontsize{8}{10}\selectfont
  \color{black!55}
  邹永康
}

\fancyfoot[R]{%
  \fontsize{8}{10}\selectfont
  \color{black!55}
  \thepage\ / \pageref*{LastPage}
}

\begin{document}

\BodyFont

% =========================================================
% 顶部：左侧姓名与职位，右侧联系方式
% 保留原有字号；较长的求职方向在左栏内自然换行
% =========================================================

\begin{minipage}[t]{0.60\textwidth}
  \vspace{0pt}

  {\zihao{-1}\heiti 邹永康\par}

  \vspace{5pt}

  {\zihao{-4}\heiti
    求职方向：Agent Harness 工程师\par
  }

  \vspace{4pt}

  {\DetailFont
    Agent Runtime · Context Engineering · Agent Evaluation\par
  }
\end{minipage}\hfill
\begin{minipage}[t]{0.38\textwidth}
  \vspace{0pt}
  \raggedleft
  \zihao{-5}

  \href{mailto:yongkang.zou.ai@gmail.com}
       {yongkang.zou.ai@gmail.com}\\
  \ResumeUrl{https://github.com/inin-zou}
            {github.com/inin-zou}\\
  \ResumeUrl{https://yongkang.dev/}
            {yongkang.dev}
\end{minipage}

\par\vspace{4pt}

% =========================================================
% 工作经历
% =========================================================

\section{工作经历}

\ResumeItem
  {Epiminds（斯德哥尔摩）}
  [Founding AI 工程师]
  [2026.02 -- 2026.05]

\begin{resumeBullets}
  \item 构建企业级\textbf{多智能体营销平台}，
  设计并实现 \textbf{20+ 个专家 Agent}，
  覆盖受众分析、广告创意、优化与跨平台投放；
  基于 \textbf{Vercel AI SDK}
  实现模型调用、结构化输出与流式交互，
  与Google、Meta Ads等平台API集成。

  \item 基于 \textbf{Temporal}
  构建长生命周期 Agent Workflow，
  将多步骤推理、工具调用、人工审批与外部 API 操作
  编排为可持久化执行流程，
  支持\textbf{状态恢复、失败重试与故障恢复}；
  负责 React、Express、PostgreSQL 与 Redis 核心系统开发
  并部署于 GCP，服务 \textbf{100+ 企业日活用户}。

  \item 设计并落地\textbf{用户反馈驱动的 Agent 持续优化闭环}：
  通过产品侧点赞 / 点踩收集低摩擦用户反馈，
  每周对生产 Trace 进行离线分析，
  从端到端回答进一步归因至
  \textbf{Agent Routing、Tool Call、Context} 等环节；
  结合用户历史 Profile 与跨 Agent 交互进行归因，
  而非孤立评估单条 Trace，
  以识别长期偏好与系统性失败模式。

  \item 参考 \textbf{SkillOpt} 思路迭代 Few-shot、Prompt / Skill
  与路由策略；基于脱敏生产数据构建\textbf{分层回归测试集}，
  在 Harness / 模型更新后执行端到端与工具级测试，
  联合评估\textbf{回答质量、延迟与执行步数}，检测行为退化。

\end{resumeBullets}

\vspace{\EntryGap}

\ResumeItem
  {Mozart AI（伦敦 / 巴黎）}
  [AI 开发工程师]
  [2025.10 -- 2026.01]

\begin{resumeBullets}
  \item 从 0 到 1 交付 \textbf{AI Music Video} 服务，
  负责架构设计、多模态生成流水线与生产部署；
  将歌词和音乐特征转化为 LLM 生成的叙事、视觉风格与分镜，
  协调下游图像与视频模型，并实现音乐节拍对齐；
  协作构建长耗时生成任务的异步执行机制并部署于 GCP，
  服务约 \textbf{700 DAU}。

  \item 参与歌词生成小模型后训练，围绕团队共同维护的歌曲样本进行脱敏与重写，
  从歌词内容、结构、韵脚等维度设计质量标准，
  并结合 \textbf{LLM Judge 与声学特征分析}
  构建小规模、高质量训练与偏好数据集。

  \item 构建 \textbf{RLAIF-style Preference Data Pipeline}：
  首先通过 \textbf{SFT} 使模型学习作曲结构框架与分段能力；
  随后采用 \textbf{DPO} 学习团队定义的高质量歌词偏好，
  保持结构正确性的同时优化 Creative Writing 能力。
\end{resumeBullets}

\vspace{\EntryGap}

\ResumeItem
  {Misogi Labs（巴黎）}
  [AI 开发工程师]
  [2025.05 -- 2025.09]

\begin{resumeBullets}
  \item 面向候选分子筛选场景，
  集成 \textbf{MCP、LangGraph 与生成式模型}，
  构建包含分子特性预测与 Tool Calling 的多智能体工作流，
  并负责端到端流程编排及计算化学场景验证。
\end{resumeBullets}

\vspace{\EntryGap}

\ResumeItem
  {法国兴业银行（巴黎）}
  [AI 开发工程师，实习]
  [2024.04 -- 2024.10]

\begin{resumeBullets}
  \item 主导开发跨部门 \textbf{Agentic RAG} 知识问答系统，
  使用 \textbf{LangChain 和 ElasticSearch} 构建文档检索，
  接入 \textbf{4,000+ 份企业内部文档}；
  同期与团队交付 CI/CD 部署自动化系统及生产事故检测助手。
\end{resumeBullets}

% =========================================================
% 精选项目
% =========================================================

\section{精选项目}

\ResumeItem
  {Codex Privacy HUD}
  [OpenAI Privacy Hackathon - Paris 冠军 · GitHub 累积 100+ Stars]
  [2026.09 -- 至今]

\begin{resumeBullets}
  \item 开发面向 \textbf{OpenAI Codex} 的运行时隐私与审计插件，
  集成本地 \textbf{openai/privacy-filter}；
  基于 \textbf{Lifecycle Hooks}
  捕获用户输入、工具结果、子智能体及工具调用参数等
  已接入边界中的敏感信息流。

  \item 构建\textbf{事件驱动的 Session Disclosure Ledger}，
  区分本地检测、已披露与已阻止事件；将 Privacy HUD 集成至独立构建的
  \textbf{Codex CLI 状态栏}，通过实时状态、会话审计及逐事件详情
  展示 Agent 执行过程中敏感数据的流转状态；
  \textbf{获奖后持续迭代，目前由本人独立维护，
  GitHub 累积 100+ Stars}。
\end{resumeBullets}

\ProjectLink{https://github.com/inin-zou/codex-privacy-hud}

% =========================================================
% 技术项目
% =========================================================

\section{技术项目}

\ResumeItem
  {Clio：全双工语音 Agent}
  [Big Berlin Hack · Inca Track 冠军]
  [2026]

\begin{resumeBullets}
  \item 基于 \textbf{PersonaPlex 7B}
  构建面向保险理赔场景的全双工 Speech-to-Speech Agent，
  通过 \textbf{Twilio / LiveKit} 接入真实电话链路；
  将实时音频执行与异步业务控制分离，
  使控制通道仅处理结构化指令与转录事件。

  \item 将信息抽取与业务推理移出实时语音 Hot Path，
  通过规则式 \textbf{Intervention Gate}
  与逐 Token 注入控制复述确认、缺失字段追问及结束检查；
  针对高风险字段加入用户原话锚定校验，过滤缺乏来源依据的状态更新。

  \item 将模型与 LiveKit Agent 部署于 A100，
  结合会话状态恢复降低启动开销；
  对电话链路、WebRTC、GPU Inference 进行 Latency Profiling，
  项目记录的端到端
  \textbf{Round-trip Latency 约为 450--650\,ms}。
\end{resumeBullets}

\ProjectLink{https://github.com/inin-zou/Clio}

\vspace{\EntryGap}

\ResumeItem
  {KernelGen 多后端 AI 算子优化}
  [GOSIM KernelGen 2026 · Sparse Attention 赛道第 1 名]
  [2026.05]

\begin{resumeBullets}
  \item 基于 \textbf{Triton / FlagTree}
  优化动态稀疏注意力与 DeepSeek mHC，
  为 Ascend、天数智芯、海光、摩尔线程与沐曦
  五类 AI 加速器设计后端分派及兼容执行路径。

  \item 针对 Top-k Gather Sparse Attention，
  结合在线 Softmax、BF16 计算路径、
  FP32 回退与 V 分块策略优化访存及寄存器压力；
  相较 PyTorch 参考实现，
  五后端平均加速 \textbf{1.97 倍}。

  \item 针对 mHC 使用运行时 Sinkhorn 循环、
  自动调参与按输入规模分派，
  在五后端取得平均 \textbf{71.85 倍加速}，
  摩尔线程后端达到 \textbf{201.32 倍}；
  建立参考实现、数值正确性测试与性能 Benchmark。
\end{resumeBullets}

\ProjectLink{https://github.com/inin-zou/kernelgen-challenge}

% =========================================================
% 技能与荣誉
% =========================================================

\section{技能与荣誉}

\begin{itemize}[
  leftmargin=0pt,
  label={},
  itemsep=4pt,
  parsep=0pt,
  topsep=0pt,
  partopsep=0pt
]
  \item \textbf{智能体开发：}
  MCP、RAG；Vercel AI SDK、LangGraph、LangChain

  \item \textbf{监控与可观测性：}
  Langfuse、Sentry、Prometheus、Grafana、PostHog

  \item \textbf{模型训练与推理：}
  SFT（LoRA）、RL (DPO、GRPO)；vLLM、Triton、PyTorch

  \item \textbf{编程语言：}
  Python、TypeScript、SQL

  \item \textbf{后端与系统工程：}
  Prisma、Node.js、Express、PostgreSQL、Redis、
  Temporal、Docker、GitHub Actions、GCP

  \item \textbf{前端开发：}
  React、Next.js、Vite、tailwind CSS

  \item \textbf{竞赛荣誉：}
  14 次国际 AI Hackathon 获奖，包括 OpenAI Hackathon 冠军、
  Lovable x Google DeepMind 黑客松季军

  \item \textbf{语言能力：}
  法语（DALF C1）、英语（IELTS 7.0）
\end{itemize}

% =========================================================
% 教育经历
% 保留学校、学位、成绩三行布局及原有字号、斜体
% =========================================================

\section{教育经历}

\Education
  {巴黎文理研究大学 -- PSL（Université Paris Dauphine -- PSL）}
  {2023.09 -- 2024.12}
  {\textnormal{计算机科学 | }硕士学位}
  {\textbf{QS World University Rankings 2025：\#24}
   | \textbf{年级排名第 10}}

\vspace{6pt}

\Education
  {巴黎萨克雷大学（Université Paris-Saclay）}
  {2022.09 -- 2023.08}
  {\textnormal{计算机科学 | }硕士一年级}
  {\textbf{QS World University Rankings 2025：\#73}
   | \textbf{GPA：14.41/20}
   | \textbf{年级排名第 3}}

\vspace{6pt}
\vspace{6pt}

\Education
  {图卢兹第一大学（Université Toulouse 1 Capitole）}
  {2019.09 -- 2022.06}
  {\textnormal{计算机、经济管理学 | }学士学位}
  {\textbf{经济学与计量经济学 QS Rankings：\#71}
   | \textbf{年级排名第 5}}

\end{document}
